\documentclass[11pt,a4paper]{article}
\usepackage{cmap}
\usepackage[T2A]{fontenc}
\usepackage[utf8]{inputenc}
\usepackage[russian]{babel}
\usepackage{amsmath,amssymb,amsthm}
\usepackage{geometry}
\geometry{left=20mm,right=20mm,top=20mm,bottom=22mm}
\usepackage{tikz}
\usetikzlibrary{arrows.meta,positioning,shapes.geometric}
\usepackage{xcolor}
\usepackage{tabularx,array,booktabs,multicol}
\usepackage{fancyhdr}
\usepackage{enumitem}
\usepackage{microtype}
\usepackage{needspace}
\usepackage{setspace}
\usepackage{hyperref}
\hypersetup{hidelinks,pdftitle={Количество делителей натурального числа},pdfauthor={Учебно-методическое пособие}}
\definecolor{ink}{HTML}{253644}
\definecolor{navy}{HTML}{324F69}
\definecolor{slate}{HTML}{798E9F}
\definecolor{muted}{HTML}{586B78}
\definecolor{pale}{HTML}{EDF2F5}
\color{ink}
\setstretch{1.08}
\setlength{\parindent}{0pt}
\setlength{\parskip}{4pt}
\setlist[itemize]{leftmargin=6mm,itemsep=2pt,topsep=3pt}
\setcounter{secnumdepth}{0}
\pagestyle{fancy}
\fancyhf{}
\fancyhead[L]{\footnotesize\color{muted} ЕГЭ по математике\,|\,Целые числа}
\fancyhead[R]{\footnotesize\color{muted} Количество делителей}
\fancyfoot[C]{\small\color{muted}\thepage}
\renewcommand{\headrulewidth}{0.35pt}
\setlength{\headheight}{14pt}
\renewcommand{\arraystretch}{1.34}
\newcommand{\taun}{\tau}
\newcommand{\blocktitle}[2]{\clearpage\section{Блок #1. #2}\vspace{-2mm}\noindent\color{slate}\rule{\linewidth}{.6pt}\color{ink}\vspace{2mm}}
\newcommand{\concept}[1]{\par\needspace{4\baselineskip}{\color{navy}\bfseries #1.}\enspace}
\newcommand{\example}[1]{\par\needspace{6\baselineskip}{\color{navy}\bfseries Разобранный пример.} #1\par}
\newcommand{\warning}[1]{\par\needspace{4\baselineskip}{\color{muted}\bfseries Типичная ошибка.} #1\par}
\newcommand{\selfcheck}[1]{\par\needspace{3\baselineskip}{\color{muted}\bfseries Проверка себя.} #1\par}
\newcommand{\minitrain}{\par\needspace{7\baselineskip}\vspace{1mm}{\color{navy}\bfseries Три мини-упражнения}\par\vspace{-1mm}}
\newcommand{\mini}[2]{\par\noindent\hspace*{2mm}\textbf{#1)}\enspace #2\par}
\newcommand{\task}[2]{\par\needspace{2\baselineskip}\noindent\makebox[8mm][l]{\color{navy}\bfseries #1.}#2\par\vspace{1.8mm}}
\newcommand{\training}[2]{\clearpage\section*{Тренировочный блок #1}\textcolor{muted}{#2}\par\vspace{3mm}\rule{\linewidth}{.4pt}\vspace{2mm}}
\newcommand{\answers}[2]{\clearpage\section*{Ответы к тренировочному блоку #1}\vspace{3mm}#2}
\newcommand{\anstable}[1]{\begin{center}\small\begin{tabularx}{\linewidth}{>{\centering\arraybackslash}p{9mm}X>{\centering\arraybackslash}p{9mm}X}\toprule\textbf{№}&\textbf{Ответ}&\textbf{№}&\textbf{Ответ}\\\midrule #1\\\bottomrule\end{tabularx}\end{center}}
\tikzset{
 st/.style={ellipse,draw=navy,line width=.8pt,fill=pale,align=center,minimum height=10mm,text width=40mm,font=\small},
 pr/.style={rectangle,rounded corners=2mm,draw=slate,line width=.8pt,align=center,minimum height=11mm,text width=65mm,font=\small},
 de/.style={diamond,aspect=2.4,draw=navy,line width=.8pt,align=center,text width=28mm,inner sep=1.7mm,font=\small},
 ar/.style={-{Stealth[length=2.4mm,width=1.7mm]},line width=.8pt,draw=navy},
 ann/.style={font=\footnotesize\bfseries,inner sep=1.5pt,color=navy}
}
\newcommand{\schemetitle}[1]{\clearpage\section*{Алгоритм: #1}\vspace{10mm}}
\begin{document}
\begin{titlepage}
\thispagestyle{empty}
\vspace*{26mm}
{\color{navy}\large УЧЕБНО-МЕТОДИЧЕСКОЕ ПОСОБИЕ}\par
\vspace{12mm}
{\fontsize{26}{34}\selectfont\bfseries Количество делителей\\[3mm] натурального числа}\par
\vspace{10mm}
{\Large Теория чисел. Подготовка к профильному ЕГЭ\par}
\vspace{19mm}
{\color{muted}\large Каноническое разложение, формула $\tau(n)$,\\[2mm] обратные и доказательные задачи}\par
\vspace{20mm}
\begin{tikzpicture}[x=1cm,y=1cm]
\draw[navy,line width=.95pt] (0,0) .. controls (2.9,.9) and (6.0,-.65) .. (8.3,.14) .. controls (10.7,.75) and (12.8,.28) .. (15.6,.48);
\draw[slate,line width=.35pt] (0,-.36) -- (15.6,-.36);
\end{tikzpicture}
\vfill
\end{titlepage}
\clearpage
\section*{Карта темы}
\textbf{Главная идея:} каждый положительный делитель натурального числа соответствует единственному выбору показателей простых множителей. Поэтому число делителей можно вычислить, не перечисляя их.
\begin{center}
\begin{tikzpicture}[node distance=7.5mm]
\node[pr,text width=111mm] (a) {\textbf{1. Находим структуру}\quad от простых множителей к формуле $\tau(n)$};
\node[pr,text width=111mm,below=of a] (b) {\textbf{2. Учитываем условия}\quad чётные, кратные заданному числу, квадраты};
\node[pr,text width=111mm,below=of b] (c) {\textbf{3. Решаем обратные задачи}\quad восстанавливаем вид числа по $\tau(n)$};
\node[pr,text width=111mm,below=of c] (d) {\textbf{4. Доказываем свойства}\quad пары делителей, квадраты, взаимная простота};
\draw[ar](a)--(b);\draw[ar](b)--(c);\draw[ar](c)--(d);
\end{tikzpicture}
\end{center}
\concept{Обозначения} $\mathbb N=\{1,2,3,\ldots\}$; запись $d\mid n$ означает, что $d$ делит $n$. Через $\tau(n)$ обозначим \emph{количество положительных делителей} числа $n$ (включая $1$ и само $n$).
\concept{Как работать} Изучите определения и примеры; после каждого подраздела выполните три мини-упражнения; затем пройдите схему, решите десять задач блока и сверьтесь с отдельной страницей ответов.
\concept{Важная граница} Число $1$ имеет ровно один положительный делитель: $\tau(1)=1$. Для $n>1$ мы используем каноническое разложение $n=p_1^{\alpha_1}\cdots p_k^{\alpha_k}$.

% БЛОК 1
\blocktitle{1}{От канонического разложения к формуле}
\subsection{1.1. Как устроены положительные делители}
\concept{Цель} Научиться перечислять делители не перебором чисел, а выбором степеней простых множителей.
\concept{Определение} Положительное число $d$ называют делителем числа $n$, если $n=dm$ при некотором натуральном $m$. Например, делители $12$ равны $1,2,3,4,6,12$.
\concept{Напоминание} Для натурального $n>1$ каноническое разложение имеет вид
\[
 n=p_1^{\alpha_1}\cdot p_2^{\alpha_2}\cdots p_k^{\alpha_k},
\]
где $p_i$ --- различные простые числа, а $\alpha_i$ --- натуральные показатели. По основной теореме арифметики это разложение единственно с точностью до порядка множителей.
\concept{Критерий} Число
\[
 d=p_1^{\beta_1}\cdots p_k^{\beta_k}
\]
делит $n$ тогда и только тогда, когда $0\le\beta_i\le\alpha_i$ при каждом $i$. Нулевой показатель означает, что соответствующий простой множитель отсутствует: $p^0=1$.
\example{Возьмём $n=60=2^2\cdot3\cdot5$. Для любого делителя $d=2^a3^b5^c$ необходимо и достаточно
\[
 a\in\{0,1,2\},\qquad b,c\in\{0,1\}.
\]
Так, число $30=2\cdot3\cdot5$ делит $60$, а $8=2^3$ не делит: показатель тройки у двойки здесь ни при чём --- не хватает третьей двойки.}
\concept{Сравнение} У $2^4$ делители $1,2,4,8,16$ соответствуют показателям $0,1,2,3,4$. У числа $2^4\cdot3$ каждый из них можно взять либо без множителя $3$, либо с ним: число вариантов удваивается.
\warning{Не следует считать только \emph{простые} делители. Например, у числа $12$ два различных простых делителя ($2$ и $3$), но шесть положительных делителей.}
\minitrain
\mini{1}{Запишите каноническое разложение $72$. ($\boldsymbol{2^3\cdot3^2}$.)}
\mini{2}{Делится ли $2^3\cdot3^2$ на $2^2\cdot3$? (\textbf{Да}.)}
\mini{3}{Сколько показателей может иметь двойка в делителе $2^4$? ($\boldsymbol{5}$.)}

\subsection{1.2. Теорема о количестве делителей}
\concept{Теорема} Если $n=p_1^{\alpha_1}\cdots p_k^{\alpha_k}>1$, то
\[
 \boxed{\tau(n)=(\alpha_1+1)(\alpha_2+1)\cdots(\alpha_k+1).}
\]
\concept{Доказательство} Для показателя $\beta_1$ имеется $\alpha_1+1$ вариантов: от нуля до $\alpha_1$. Для $\beta_2$ --- $\alpha_2+1$ вариантов и так далее. Выборы не зависят друг от друга, а разные наборы показателей дают разные делители. По правилу произведения получаем формулу.
\example{Найдём число положительных делителей $1440$:
\begin{align*}
1440&=144\cdot10=(2^4\cdot3^2)(2\cdot5)
=2^5\cdot3^2\cdot5,\\
\tau(1440)&=(5+1)(2+1)(1+1)=6\cdot3\cdot2=36.
\end{align*}
Мы не выписывали все делители, но учли каждый ровно один раз.}
\example{Для $n=72=2^3\cdot3^2$ выпишем все комбинации степеней:
\begin{center}
\small\begin{tabular}{c|ccc}\toprule
&$3^0$&$3^1$&$3^2$\\\midrule
$2^0$&$1$&$3$&$9$\\
$2^1$&$2$&$6$&$18$\\
$2^2$&$4$&$12$&$36$\\
$2^3$&$8$&$24$&$72$\\\bottomrule
\end{tabular}
\end{center}
В таблице $4\cdot3=12$ клеток, следовательно, $\tau(72)=12$.}
\concept{Частные случаи} У простого числа $p$ ровно два делителя. У степени $p^a$ ровно $a+1$ делитель. У единицы один делитель --- формула для пустого произведения также согласуется с $\tau(1)=1$.
\selfcheck{Быстрая контрольная проверка: количество делителей числа $2^3\cdot3^2$ равно $4\cdot3$, а не $3\cdot2$ и не $2^3\cdot3^2$.}
\minitrain
\mini{1}{Найдите $\tau(12)$. ($\boldsymbol{6}$.)}
\mini{2}{Найдите $\tau(2^3\cdot5^2)$. ($\boldsymbol{12}$.)}
\mini{3}{Найдите $\tau(1)$. ($\boldsymbol{1}$.)}

\schemetitle{вычисление количества делителей}
\begin{center}
\begin{tikzpicture}[node distance=12mm]
\node[st] (s) {Дано натуральное $n$};
\node[de,below=16mm of s] (q) {$n=1$?};
\node[pr,below=19mm of q] (a) {Разложите $n$ на простые множители};
\node[pr,below=of a] (b) {Запишите все показатели $\alpha_1,\ldots,\alpha_k$};
\node[pr,below=of b] (c) {Вычислите произведение $(\alpha_1+1)\cdots(\alpha_k+1)$};
\node[st,below=12mm of c] (f) {Ответ: $\tau(n)$};
\node[pr,text width=27mm,right=26mm of q] (one) {$\tau(1)=1$};
\draw[ar](s)--(q);
\draw[ar](q.south)--node[ann,right]{нет}(a.north);
\draw[ar](q.east)--node[ann,above]{да}(one.west);
\draw[ar](a)--(b);\draw[ar](b)--(c);\draw[ar](c)--(f);
\draw[ar](one.east)--++(9mm,0)|-(f.east);
\end{tikzpicture}
\end{center}

\training{1}{Основные вычисления: от канонической записи к значению $\tau(n)$.}
\task{1}{Вычислите $\tau(18)$.}
\task{2}{Вычислите $\tau(45)$.}
\task{3}{Сколько положительных делителей имеет число $80$?}
\task{4}{Сколько положительных делителей имеет число $150$?}
\task{5}{Вычислите $\tau(216)$.}
\task{6}{Найдите количество положительных делителей числа $840$.}
\task{7}{Найдите $\tau(2^5\cdot3^2\cdot11)$.}
\task{8}{Сколько положительных делителей у числа $1440$?}
\task{9}{Вычислите $\tau(3600)$.}
\task{10}{Сколько положительных делителей имеет число $10^6$?}
\answers{1}{\anstable{
1&$6$&6&$32$\\
2&$6$&7&$36$\\
3&$10$&8&$36$\\
4&$12$&9&$45$\\
5&$16$&10&$49$
}}
\selfcheck{Для задания 9: $3600=2^4\cdot3^2\cdot5^2$, поэтому $\tau(3600)=5\cdot3\cdot3=45$.}

% БЛОК 2
\blocktitle{2}{Подсчёт делителей с дополнительными условиями}
\subsection{2.1. Кратность, чётность и запрет делителей}
\concept{Цель} Научиться переводить условия на делитель в ограничения на показатели его простых множителей.
\concept{Основной приём} Пусть $n=2^a3^b5^c$, а $d=2^x3^y5^z$ --- делитель $n$. Тогда показатели $x,y,z$ ограничены сверху числами $a,b,c$ соответственно. Любое дополнительное требование меняет \emph{допустимые значения} этих показателей.
\concept{Сравнение трёх условий} Для $n=2^4\cdot3^2\cdot5$:
\begin{itemize}
\item условие «$d$ нечётен» задаёт $x=0$;
\item условие «$12\mid d$» задаёт $x\ge2$, $y\ge1$;
\item условие «$3\nmid d$» задаёт $y=0$.
\end{itemize}
\example{Сколько делителей числа $360=2^3\cdot3^2\cdot5$ кратны $12$?
Запишем $d=2^x3^y5^z$. Условие $12=2^2\cdot3\mid d$ означает
\[
 x\in\{2,3\},\qquad y\in\{1,2\},\qquad z\in\{0,1\}.
\]
Следовательно, подходящих делителей $2\cdot2\cdot2=8$. Проверка перечислением:
\[12,\ 24,\ 36,\ 60,\ 72,\ 120,\ 180,\ 360.\]
Получили ровно восемь чисел.}
\example{Сколько нечётных делителей у $360$? Нужно взять $x=0$; для $y$ остаются три варианта, для $z$ --- два. Получаем $3\cdot2=6$. Всего делителей $4\cdot3\cdot2=24$, значит, чётных $24-6=18$.}
\concept{Удобное сокращение} Если $a\mid n$, то количество положительных делителей $n$, кратных $a$, равно $\tau(n/a)$. Действительно, каждому делителю $d$ вида $d=ak$ соответствует ровно один положительный делитель $k\mid n/a$. Это работает именно при $a\mid n$.
\warning{Условие «кратно $12$» не означает «содержит ровно две двойки и ровно одну тройку». Показатели могут быть и больше минимально необходимых.}
\minitrain
\mini{1}{Сколько нечётных делителей у числа $72$? ($\boldsymbol{3}$.)}
\mini{2}{Сколько чётных делителей у числа $72$? ($\boldsymbol{9}$.)}
\mini{3}{Сколько делителей числа $72$ кратны $6$? ($\boldsymbol{6}$.)}

\subsection{2.2. Делители-квадраты и делители-кубы}
\concept{Ключевая идея} Натуральное число является полным квадратом тогда и только тогда, когда \emph{все} показатели в его каноническом разложении чётные. Для полного куба все показатели кратны трём.
\concept{Формула} Для $n=p_1^{\alpha_1}\cdots p_k^{\alpha_k}$ количество его делителей, являющихся квадратами, равно
\[
 \boxed{\prod_{i=1}^{k}\left(\left\lfloor\frac{\alpha_i}{2}\right\rfloor+1\right).}
\]
Здесь $\lfloor t\rfloor$ --- целая часть числа $t$ (наибольшее целое, не превосходящее $t$). Для степеней $p_i$ можно выбрать показатели $0,2,4,\ldots$, не превосходящие $\alpha_i$.
\concept{Обобщение} Количество положительных делителей, которые являются полными $r$-ми степенями ($r\in\mathbb N$), равно
\[
 \prod_{i=1}^k\left(\left\lfloor\frac{\alpha_i}{r}\right\rfloor+1\right).
\]
\example{Рассмотрим $n=2^5\cdot3^4\cdot7$. Для делителя-квадрата показатели могут быть такими:
\[
 \begin{array}{c|ccc}
 \text{простое}&2&3&7\\\hline
 \text{допустимые степени}&0,2,4&0,2,4&0
 \end{array}
\]
Всего $3\cdot3\cdot1=9$ делителей-квадратов. А делителей-кубов $\left(\lfloor5/3\rfloor+1\right)\left(\lfloor4/3\rfloor+1\right)(\lfloor1/3\rfloor+1)=2\cdot2\cdot1=4$.}
\example{Найдём количество делителей числа $2^4\cdot3^3\cdot5^2$, которые одновременно \emph{кратны $45$ и являются квадратами}. Так как $45=3^2\cdot5$, то для показателей имеем
\[
 \beta_2\in\{0,2,4\},\quad \beta_3=2,\quad \beta_5=2.
\]
Получается три делителя: $3^2\cdot5^2$, $2^2\cdot3^2\cdot5^2$, $2^4\cdot3^2\cdot5^2$.}
\warning{Квадратный делитель может равняться $1$, ведь $1=1^2$. Даже если $n$ не квадрат, некоторые его делители могут быть квадратами.}
\minitrain
\mini{1}{Сколько делителей-квадратов у $144$? ($\boldsymbol{6}$.)}
\mini{2}{Сколько делителей числа $2^5\cdot3^2$ имеют вид $3^k$? ($\boldsymbol{3}$.)}
\mini{3}{Сколько делителей-кубов у $2^4\cdot5^2$? ($\boldsymbol{2}$.)}

\schemetitle{подсчёт делителей с ограничениями}
\begin{center}
\begin{tikzpicture}[node distance=12mm]
\node[st] (s) {Даны число $n$ и условие};
\node[pr,below=of s] (a) {Разложите $n$ на простые множители};
\node[pr,below=of a] (b) {Переведите условие в допустимые показатели $\beta_i$};
\node[de,below=16mm of b] (q) {Есть допустимые показатели?};
\node[pr,below=19mm of q] (c) {Умножьте числа вариантов для каждого $\beta_i$};
\node[st,below=13mm of c] (f) {Количество делителей};
\node[pr,text width=26mm,right=25mm of q] (zero) {Ответ: $0$};
\draw[ar](s)--(a);\draw[ar](a)--(b);\draw[ar](b)--(q);
\draw[ar](q.south)--node[ann,right]{да}(c.north);
\draw[ar](q.east)--node[ann,above]{нет}(zero.west);
\draw[ar](c)--(f);
\draw[ar](zero.east)--++(8mm,0)|-(f.east);
\end{tikzpicture}
\end{center}

\training{2}{Считаем только те делители, которые удовлетворяют заданным условиям.}
\task{1}{Сколько нечётных положительных делителей у числа $360$?}
\task{2}{Сколько чётных положительных делителей у числа $360$?}
\task{3}{Сколько положительных делителей числа $180$ кратны $5$?}
\task{4}{Сколько положительных делителей числа $2^4\cdot3^3$ кратны $12$?}
\task{5}{Сколько делителей числа $900$ являются полными квадратами?}
\task{6}{Сколько делителей-квадратов у числа $2^5\cdot3^4\cdot7$?}
\task{7}{Сколько делителей числа $2^5\cdot3^2\cdot5$ кратны $6$?}
\task{8}{Сколько положительных делителей числа $720$ кратны $10$, но \emph{не} кратны $3$?}
\task{9}{Сколько положительных делителей числа $2^5\cdot3^4\cdot5^2$ кратны $18$, но не кратны $4$?}
\task{10}{Сколько делителей числа $2^4\cdot3^3\cdot5^2$ одновременно кратны $45$ и являются квадратами?}
\answers{2}{\anstable{
1&$6$&6&$9$\\
2&$18$&7&$20$\\
3&$9$&8&$4$\\
4&$9$&9&$9$\\
5&$8$&10&$3$
}}
\selfcheck{В задаче 8: $720=2^4\cdot3^2\cdot5$. Для делителя, кратного $10$, но не $3$, показатель двойки может быть $1,2,3,4$, тройки --- только $0$, пятёрки --- только $1$. Ответ: $4$.}

% БЛОК 3
\blocktitle{3}{Обратные задачи: сколько делителей должно быть}
\subsection{3.1. Что можно узнать о числе по значению $\tau(n)$}
\concept{Цель} Восстановить возможные формы канонического разложения по заданному количеству делителей.
\concept{Обратный переход} Если $\tau(n)=m$, то необходимо представить $m$ в виде произведения натуральных множителей, больших единицы:
\[
 m=(\alpha_1+1)\cdots(\alpha_k+1).
\]
После этого показатели простых множителей числа $n$ равны $\alpha_i$. Каждая \emph{неупорядоченная} факторизация $m$ задаёт один набор показателей, но сами простые числа можно выбирать по-разному.
\concept{Сравнение} Если $\tau(n)=3$, то $3$ --- простое, а значит, $n=p^2$. При $\tau(n)=4$ есть два варианта: $4=4$ даёт $n=p^3$, а $4=2\cdot2$ даёт $n=pq$ при различных простых $p,q$.
\example{Найдём все \emph{возможные виды} натуральных чисел с шестью делителями:
\[
 6=6=3\cdot2.
\]
Отсюда $n=p^5$ либо $n=p^2q$, где $p,q$ --- различные простые числа. Например, $32=2^5$ и $12=2^2\cdot3$ имеют по шесть делителей. Иных видов нет: иных разложений числа $6$ на множители, большие единицы, не существует.}
\concept{Почему простое $m$ особенно удобно} Если $\tau(n)=m$, где $m$ простое, то обязательно $n=p^{m-1}$. Например, при $\tau(n)=11$ любое такое число является десятой степенью простого.
\warning{Не путайте утверждение «$\tau(n)=6$» с «$n$ имеет шесть \emph{простых} делителей». Количество всех делителей и число различных простых множителей --- разные характеристики.}
\minitrain
\mini{1}{Если $\tau(p^a)=5$, найдите $a$. ($\boldsymbol{4}$.)}
\mini{2}{Каков вид числа $n$, если $\tau(n)=3$? ($\boldsymbol{p^2}$, где $p$ простое.)}
\mini{3}{Назовите все виды $n$ при $\tau(n)=4$. ($\boldsymbol{p^3}$ или $\boldsymbol{pq}$, $p\ne q$ --- простые.)}

\subsection{3.2. Наименьшее число с заданным количеством делителей}
\concept{Цель} Найти \emph{наименьшее} натуральное число $n$ при условии $\tau(n)=m$.
\concept{Принцип минимальности} При заданных показателях $\alpha_i$ число будет минимальным, если использовать \emph{наименьшие простые} $2,3,5,\ldots$; большие показатели ставятся при меньших простых. В самом деле, если $p<q$ и $a>b$, то $p^aq^b<p^bq^a$, поскольку $(q/p)^{a-b}>1$.
\concept{Ограничение} Этот принцип выбирает минимум \emph{для фиксированного набора показателей}. Надо ещё сравнить разные разложения числа $m$ на множители $\alpha_i+1$.
\example{Найдём наименьшее число с $12$ делителями. Все неупорядоченные разложения числа $12$ на множители, большие единицы:
\[
 12;\qquad6\cdot2;\qquad4\cdot3;\qquad3\cdot2\cdot2.
\]
Соответствующие минимальные числа:
\[
 \begin{array}{c|c|r}
 \text{множители }12&\text{число}&\text{значение}\\\hline
 12&2^{11}&2048\\
 6\cdot2&2^5\cdot3&96\\
 4\cdot3&2^3\cdot3^2&72\\
 3\cdot2\cdot2&2^2\cdot3\cdot5&60
 \end{array}
\]
Итак, минимум равен $\boxed{60}$, причём $\tau(60)=3\cdot2\cdot2=12$.}
\example{Найдём минимальное число, имеющее ровно девять делителей. Возможны формы $p^8$ и $p^2q^2$. Минимумы: $2^8=256$ и $2^2\cdot3^2=36$. Следовательно, ответ $36$.}
\selfcheck{Для доказательства \emph{наименьшего} числа недостаточно привести один подходящий пример: нужно рассмотреть все допустимые наборы показателей и сравнить их минимальные реализации.}
\minitrain
\mini{1}{Какое наименьшее число имеет $5$ делителей? ($\boldsymbol{16}$.)}
\mini{2}{Какое наименьшее число имеет $6$ делителей? ($\boldsymbol{12}$.)}
\mini{3}{Какое наименьшее число имеет $10$ делителей? ($\boldsymbol{48}$.)}

\schemetitle{поиск наименьшего числа с $m$ делителями}
\begin{center}
\begin{tikzpicture}[node distance=12mm]
\node[st] (s) {Задано $m=\tau(n)$};
\node[pr,below=of s] (a) {Перечислите факторизации $m$ на множители $\ge2$};
\node[pr,below=of a] (b) {Возьмите очередной набор; вычтите $1$ из его множителей};
\node[pr,below=of b] (c) {Поставьте большие показатели при меньших простых};
\node[pr,below=of c] (d) {Найдите наименьшее число для каждого набора};
\node[de,below=17mm of d] (q) {Все наборы проверены?};
\node[st,below=17mm of q] (f) {Выберите наименьшее};
\draw[ar](s)--(a);\draw[ar](a)--(b);\draw[ar](b)--(c);\draw[ar](c)--(d);\draw[ar](d)--(q);
\draw[ar](q.south)--node[ann,right]{да}(f.north);
\draw[ar](q.west)--node[ann,above]{нет}++(-17mm,0)|-([yshift=-2mm]b.west);
\end{tikzpicture}
\end{center}

\training{3}{Обратная формула, виды чисел с заданным $\tau(n)$ и доказательство минимальности.}
\task{1}{Какой вид имеет натуральное число $n$, если $\tau(n)=2$?}
\task{2}{Найдите $a$, если $\tau(3^a)=7$ и $a\in\mathbb N$.}
\task{3}{Найдите наименьшее натуральное число с тремя положительными делителями.}
\task{4}{Найдите наименьшее натуральное число с четырьмя положительными делителями.}
\task{5}{Найдите наименьшее натуральное число с шестью положительными делителями.}
\task{6}{Укажите все возможные виды натуральных чисел с шестью делителями (через различные простые $p,q$).}
\task{7}{Найдите наименьшее натуральное число с восемью положительными делителями.}
\task{8}{Найдите наименьшее натуральное число с девятью положительными делителями.}
\task{9}{Найдите наименьшее натуральное число с двенадцатью положительными делителями.}
\task{10}{Найдите наименьшее натуральное число, которое делится на $18$ и имеет ровно $12$ положительных делителей.}
\answers{3}{\anstable{
1&$n=p$, $p$ --- простое&6&$p^5$ или $p^2q$, $p\ne q$\\
2&$a=6$&7&$24$\\
3&$4$&8&$36$\\
4&$6$&9&$60$\\
5&$12$&10&$72$
}}
\selfcheck{В задании 10: $72=2^3\cdot3^2$ кратно $18$ и имеет $(3+1)(2+1)=12$ делителей. Числа меньше $72$, кратные $18$, --- $18,36,54$; число их делителей соответственно $6,9,8$.}

% БЛОК 4
\blocktitle{4}{Доказательства и задачи повышенного уровня}
\subsection{4.1. Парность делителей и полные квадраты}
\concept{Наблюдение} Делители числа $n$ удобно объединять в пары $d$ и $n/d$. Например, у числа $30$ пары
\[
 (1,30),\ (2,15),\ (3,10),\ (5,6).
\]
У числа $36$ получаем $(1,36)$, $(2,18)$, $(3,12)$, $(4,9)$ и особый делитель $6$, для которого $6\cdot6=36$.
\concept{Теорема} Число положительных делителей натурального числа $n$ \emph{нечётно тогда и только тогда}, когда $n$ --- полный квадрат.
\concept{Доказательство через пары} Если $n$ не квадрат, равенство $d=n/d$ невозможно, поэтому все делители разбиваются на пары разных чисел и их количество чётно. Если $n=s^2$, то у делителя $s$ пара совпадает с ним самим. Все остальные делители разбиваются на пары; значит, количество делителей нечётно.
\concept{Второе доказательство (через формулу)} Если $n=p_1^{\alpha_1}\cdots p_k^{\alpha_k}$, то произведение $(\alpha_1+1)\cdots(\alpha_k+1)$ нечётно тогда и только тогда, когда все $\alpha_i$ чётны. Но именно так записывается полный квадрат.
\example{Пусть $n=144=2^4\cdot3^2$. Тогда
\[
 \tau(144)=(4+1)(2+1)=15.
\]
Из $15$ делителей семь пар имеют разные элементы, а делитель $12=\sqrt{144}$ остаётся без второго, отличного от него, элемента. В результате $7\cdot2+1=15$.}
\concept{Следствие} Если $\tau(n)=2k$ и $n$ не квадрат, то существует ровно $k$ пар $(d,n/d)$ с $d<n/d$. Если $\tau(n)=2k+1$, то $n$ квадрат, имеются $k$ различных пар и центральный делитель $\sqrt n$.
\warning{Нечётность числа делителей не означает, что число само нечётно: $36$ --- чётный квадрат, но $\tau(36)=9$.}
\minitrain
\mini{1}{Найдите $\tau(25)$. ($\boldsymbol{3}$.)}
\mini{2}{Найдите $\tau(30)$. ($\boldsymbol{8}$.)}
\mini{3}{Является ли $2^2\cdot3^4$ квадратом? (\textbf{Да}; $\tau=15$.)}

\subsection{4.2. Произведение взаимно простых чисел и конструкции}
\concept{Теорема о взаимно простых множителях} Если $\text{НОД}(a,b)=1$, то
\[
 \boxed{\tau(ab)=\tau(a)\tau(b).}
\]
\concept{Доказательство} При $\text{НОД}(a,b)=1$ в канонических разложениях $a$ и $b$ нет общих простых множителей. Поэтому показатели этих множителей в $ab$ просто объединяются, а множители в формуле для $\tau$ образуют произведение $\tau(a)\tau(b)$.
\concept{Важное сравнение} Если $a=8=2^3$, $b=25=5^2$, то
\[
 \tau(200)=\tau(2^3\cdot5^2)=4\cdot3=12=\tau(8)\tau(25).
\]
Если же $a=12=2^2\cdot3$, $b=18=2\cdot3^2$, то
\[
 \tau(216)=4\cdot4=16<\tau(12)\tau(18)=6\cdot6=36.
\]
Общий простой множитель меняет показатель с $u,v$ на $u+v$, и $u+v+1<(u+1)(v+1)$ при $u,v\ge1$.
\concept{Следствие} Для любых натуральных $a,b$ выполняется $\tau(ab)\le\tau(a)\tau(b)$, причём равенство достигается тогда и только тогда, когда $a,b$ взаимно просты. Это включает случай $a=1$ или $b=1$.
\example{Найдём наименьший \emph{полный квадрат} с $45$ делителями. Поскольку число квадрат, все показатели в его каноническом разложении чётные, а значит, все множители $\alpha_i+1$ нечётные. Возможные разложения числа $45$ на множители, большие единицы:
\[
 45;\quad15\cdot3;\quad9\cdot5;\quad5\cdot3\cdot3.
\]
Сравним минимальные реализации:
\[
\begin{array}{c|c|r}
45&2^{44}&\text{больше }3600\\
15\cdot3&2^{14}\cdot3^2&147456\\
9\cdot5&2^8\cdot3^4&20736\\
5\cdot3\cdot3&2^4\cdot3^2\cdot5^2&3600
\end{array}
\]
Искомое число $\boxed{3600=60^2}$; действительно, $\tau(3600)=5\cdot3\cdot3=45$.}
\concept{Перенос идеи} Для любого заданного $m\ge2$ существует натуральное число ровно с $m$ делителями: достаточно взять $n=2^{m-1}$. Однако \emph{минимальным} такое число, вообще говоря, не будет.
\warning{Формулу $\tau(ab)=\tau(a)\tau(b)$ нельзя применять без проверки взаимной простоты. Для $a=b=2$ она предсказывала бы $4$ вместо $\tau(4)=3$.}
\minitrain
\mini{1}{Найдите $\tau(8\cdot25)$. ($\boldsymbol{12}$.)}
\mini{2}{Найдите $\tau(12\cdot18)$. ($\boldsymbol{16}$.)}
\mini{3}{Приведите пример квадрата с $21$ делителем. ($\boldsymbol{576=2^6\cdot3^2}$.)}

\schemetitle{определение чётности числа делителей}
\begin{center}
\begin{tikzpicture}[node distance=12mm]
\node[st] (s) {Дано натуральное $n$};
\node[pr,below=of s] (a) {Получите каноническое разложение $n$};
\node[de,below=17mm of a] (q) {Все показатели чётные?};
\node[pr,text width=48mm,below left=24mm and 18mm of q] (yes) {$n$ --- полный квадрат; $\tau(n)$ нечётно};
\node[pr,text width=48mm,below right=24mm and 18mm of q] (no) {$n$ не квадрат; $\tau(n)$ чётно};
\node[st,below=34mm of q] (f) {Вывод доказан};
\draw[ar](s)--(a);\draw[ar](a)--(q);
\draw[ar](q.south west)--node[ann,left]{да}(yes.north);
\draw[ar](q.south east)--node[ann,right]{нет}(no.north);
\draw[ar](yes.south)|-(f.west);
\draw[ar](no.south)|-(f.east);
\end{tikzpicture}
\end{center}

\training{4}{Парность количества делителей, взаимная простота и нестандартные конструкции.}
\task{1}{Сколько положительных делителей у числа $49$? Чётно ли это количество?}
\task{2}{Сколько положительных делителей у числа $72$? Чётно ли это количество?}
\task{3}{Является ли число $2^4\cdot3^2\cdot5^6$ полным квадратом? Чему равно количество его положительных делителей?}
\task{4}{Докажите: если число положительных делителей $n$ нечётно, то $n$ --- полный квадрат.}
\task{5}{На сколько пар различных взаимодополнительных делителей $(d,n/d)$ разбиваются делители числа $108$?}
\task{6}{Сколько пар \emph{разных} взаимодополнительных делителей у числа $144$? Сколько делителей не входит ни в одну такую пару?}
\task{7}{Вычислите $\tau(18\cdot25)$ и проверьте формулу произведения для взаимно простых множителей.}
\task{8}{Сравните $\tau(36\cdot10)$ и $\tau(36)\tau(10)$. Запишите верный знак сравнения.}
\task{9}{Найдите все натуральные числа вида $2^a3^b$, где $a,b\ge1$ --- целые и $\tau(2^a3^b)=15$.}
\task{10}{Найдите наименьший полный квадрат, имеющий ровно $45$ положительных делителей. Обоснуйте минимальность.}
\answers{4}{\anstable{
1&$3$, нечётно&6&$7$ пар; ещё $1$ делитель ($12$)\\
2&$12$, чётно&7&$18=6\cdot3$\\
3&Да; $5\cdot3\cdot7=105$&8&$24<36$\\
4&Все показатели простых чётные&9&$144$ и $324$\\
5&$6$ пар&10&$3600$
}}
\selfcheck{Для задачи 9: $(a+1)(b+1)=15=3\cdot5$, поэтому $(a,b)=(2,4)$ или $(4,2)$. Получаем $2^23^4=324$ и $2^43^2=144$.}

\clearpage
\section*{Итоговая обзорная контрольная работа}
\textcolor{muted}{10 заданий: 1--3 --- базовый уровень, 4--7 --- средний, 8--9 --- повышенный, 10 --- комплексная задача.}
\par\vspace{3mm}\rule{\linewidth}{.4pt}\vspace{2mm}
\task{1}{Вычислите $\tau(98)$.}
\task{2}{Сколько положительных делителей имеет число $256$?}
\task{3}{Найдите $\tau(2^3\cdot3^2\cdot5)$.}
\task{4}{Сколько нечётных положительных делителей имеет число $600$?}
\task{5}{Сколько положительных делителей-квадратов имеет число $1296$?}
\task{6}{Сколько положительных делителей числа $2^4\cdot3^2\cdot7$ кратны $6$?}
\task{7}{Найдите наименьшее натуральное число с $15$ положительными делителями.}
\task{8}{Докажите, что всякое натуральное число с ровно $11$ положительными делителями является полным квадратом.}
\task{9}{Сколько делителей числа $2^6\cdot3^4\cdot5^2$ одновременно кратны $90$ и являются полными квадратами?}
\task{10}{Найдите \emph{наименьшее} натуральное число, кратное $20$ и имеющее ровно $18$ положительных делителей. Докажите, что меньшего числа не существует.}
\clearpage
\section*{Ответы к итоговой обзорной контрольной работе}
\anstable{
1&$6$&6&$16$\\
2&$9$&7&$144$\\
3&$24$&8&$n=p^{10}$, значит, квадрат\\
4&$6$&9&$6$\\
5&$9$&10&$180$
}
\concept{Критерий для задач с доказательством} В № 8 из простоты числа $11$ следует, что $\tau(n)=11$ может получиться лишь при одном простом множителе и показателе $10$. В № 10 необходимо сравнить все допустимые наборы показателей, совместимые с делимостью на $20=2^2\cdot5$.
\concept{Проверка задачи 10} Число $180=2^2\cdot3^2\cdot5$ кратно $20$ и имеет $3\cdot3\cdot2=18$ делителей. Минимальность удобно доказать перебором вариантов $n<180$, кратных $20$: это $20,40,60,80,100,120,140,160$. Их количества делителей соответственно $6,8,12,10,9,16,12,12$, и ни одно не равно $18$.

\clearpage
\section*{Памятка: количество делителей натурального числа}
\concept{Главная формула} Для $n=p_1^{\alpha_1}\cdots p_k^{\alpha_k}>1$:
\[
 \boxed{\tau(n)=\prod_{i=1}^k(\alpha_i+1)},\qquad \boxed{\tau(1)=1}.
\]
\concept{Как выглядит делитель} $d=p_1^{\beta_1}\cdots p_k^{\beta_k}$, где $0\le\beta_i\le\alpha_i$.
\concept{Делители, кратные $a$} Если $a\mid n$, то их количество равно $\tau(n/a)$; если $a\nmid n$, то их нет.
\concept{Делители-квадраты и делители-кубы}
\[
 \tau_{\text{кв}}(n)=\prod_{i=1}^k\left(\left\lfloor\frac{\alpha_i}{2}\right\rfloor+1\right),\quad
 \tau_{\text{куб}}(n)=\prod_{i=1}^k\left(\left\lfloor\frac{\alpha_i}{3}\right\rfloor+1\right).
\]
\concept{Чётность} $\tau(n)$ нечётно $\Longleftrightarrow$ $n$ --- полный квадрат.
\concept{Взаимная простота} Если $\text{НОД}(a,b)=1$, то $\tau(ab)=\tau(a)\tau(b)$. Без этого условия равенство вообще говоря неверно; всегда $\tau(ab)\le\tau(a)\tau(b)$.
\concept{Обратные задачи} Для $\tau(n)=m$ найдите все разложения $m=(\alpha_1+1)\cdots(\alpha_k+1)$. Чтобы получить минимум для выбранного набора показателей, ставьте наибольший показатель при $2$, следующий при $3$ и так далее; затем сравните все наборы.
\concept{Самопроверка} Обязательно учитывайте $1$ и $n$. Для задачи на минимум обосновывайте, что все варианты охвачены. Для условий «не делится», «кратно», «квадрат» переводите слова в диапазоны и свойства показателей.
\vfill
\noindent\color{slate}\rule{\linewidth}{.45pt}\color{ink}
\par\smallskip
\textcolor{muted}{Основной навык: вместо перебора самих делителей систематически считать допустимые комбинации простых степеней.}
\end{document}